%!TEX program = lualatex
\documentclass[a5paper,ngerman,headheight=80pt,10pt,DIV=calc]{scrartcl}
\usepackage[ngerman]{babel}
\usepackage{blindtext}
\usepackage[utf8]{luainputenc}
\usepackage{scrlayer-scrpage}
\usepackage[T1]{fontenc}
\usepackage{fontspec}
\usepackage{multicol}
\usepackage{enumitem}
\usepackage{epigraph}
\usepackage[german=quotes]{csquotes}
\usepackage[footskip=20pt,bottom=50pt]{geometry}
\setmainfont{Source Sans Pro}
\MakeOuterQuote{"}

\setlength{\parindent}{0pt}
\setlength{\parskip}{4pt}

\pagestyle{plain.scrheadings}
\thispagestyle{scrheadings}
\lohead{\vspace*{-1.1cm}\includegraphics[scale=0.3]{../images/stupa_logo_rot.pdf}}
\rohead{\\Präsidium des\\Studierendenparlaments}

\lofoot[]{V.i.S.d.P. Gunhild Berdal, Jim Martens, Ramon Weilinger\\Von-Melle-Park 5, 20146 Hamburg}
\cofoot[]{}
\cohead[\thepage]{}

\renewcommand{\thesubsection}{TOP \arabic{subsection}}
\renewcommand{\thesubsubsection}{\arabic{subsubsection}.}
\setcounter{subsection}{-1}

\begin{document}
    UHH · StuPa-Präsidium · Von-Melle-Park 5 · D-20146 Hamburg\\
    Tel.: 040/450 204 39 · stupa@uni-hamburg.de

    \subsection*{Von Brotgelehrten und philosophischen Köpfen: Fremdbestimmt oder
    selbstgestaltend studieren?}

    "Was heißt und zu welchem Ende studiert man Universalgeschichte?" war die
    Frage und der Titel der Antrittsvorlesung von Friedrich Schiller im Jahr 1789.
    Darin unterscheidet er zwischen dem "Brodgelehrten", der bemüht und brav für
    gesellschaftliche Anerkennung studiert, und dem "philosophischen Kopf", der
    für Wissenschaft und Wahrheit wirkt:

    \blockquote{
        Wie ganz anders verhält sich der
        philosophische Kopf! – Eben so sorgfältig, als der seine Wissenschaft von
        allen übrigen absondert, bestrebt sich jener, ihr Gebiet zu erweitern, und
        ihren Bund mit den übrigen wieder herzustellen – herzustellen, sage ich,
        denn nur der abstrahirende Verstand hat jene Grenzen gemacht, hat jene
        Wissenschaften von einander geschieden. Wo der Brodgelehrte trennt,
        vereinigt der philosophische Geist.
    }

    Ist dieser Gegensatz noch aktuell? Führt(e) die Bologna-Reform zur
    Herausbildung von Fachidioten, die Zusammenhänge nicht begreifen und muss
    ein Universalstudium für alle die wissenschaftlichen und gesellschaftlichen
    Zusammenhänge stärker in den Fokus nehmen?

    Fremdbestimmt oder selbstgestaltend studieren? – Um im Sinne Schillers den
    Gesamtzusammenhang herzustellen, laden wir alle Interessierten zum
    philosophisch-praktischen Diskurs ein.

    Die Vizepräsidentin für Studium und Lehre, Frau Prof. Dr. Rupp, und eine
    studentische Vertreterin werden jeweils einen Input geben.

    \begin{center}
        \textbf{
        Donnerstag, 24. November 2016,\\
        18 Uhr\\
        im Raum 0079 des WiWi-Bunkers\\
        (Von-Melle-Park 5, 20146 Hamburg)}
    \end{center}
\end{document}
